\documentclass[11pt]{article}
\usepackage[margin=1in]{geometry}
\usepackage{amsmath,amssymb,amsthm}
\newtheorem{definition}{Definition}
\newtheorem{proposition}{Proposition}
\newtheorem{remark}{Remark}
\DeclareMathOperator{\tr}{tr}
\begin{document}

\section*{Step 2 Structural Statement: Complementary LQG Import Audit}

\begin{definition}[Lead complementary recognition source]
For Step 2, the lead named recognition source is loop quantum gravity /
canonical quantum geometry / spin-foam quantum geometry.  Its declared source
atoms are Ashtekar variables, spin-network states and area-volume operators,
and spin-foam transition-amplitude language.  Its declared caveats are the
semiclassical limit, dynamics and constraint control, anomaly questions, and
matter-coupling obligations.
\end{definition}

\begin{definition}[Reused finite carrier]
At refinement level \(r\in\{2,4\}\), reuse the Step 1 carrier
\[
  E_r=\mathbb{R}^{r+2},\qquad C=I.
\]
The first \(r\) coordinates are boundary entanglement-ledger coordinates
\(s_1,\ldots,s_r\).  The last two coordinates are \(g\), a
geometry-amplitude coordinate, and \(m\), a P3 route-mismatch coordinate.

The LQG-like native lens is
\[
  L^{\mathrm{LQG}}_r=\{g,\;m+0.30g\}.
\]
The surviving Step 1 target sector is
\[
  D^{\mathrm{surv}}_r=\{g,\;m+0.30g\}.
\]
The RT-like area-shadow sector is
\[
  D^{\mathrm{area}}_r=\{s_i/\sqrt r:1\le i\le r\}.
\]
\end{definition}

\begin{proposition}[Complementary finite-carrier recognition]
On the declared carrier,
\[
  \Xi_I(D^{\mathrm{surv}}_r\mid L^{\mathrm{LQG}}_r)=0
\]
up to numerical precision for \(r=2,4\).  At the same two refinement levels,
\[
  \frac{\tr\,\Xi_I(D^{\mathrm{area}}_r\mid L^{\mathrm{LQG}}_r)}
       {\tr K_{DD}^{\mathrm{area}}}=1.0 .
\]
\end{proposition}

\begin{proof}
The surviving target rows are exactly in the span of the LQG-like source rows,
so exact adequacy holds for that child sector.  The area-shadow rows live
entirely in the boundary entanglement-ledger coordinates, which are orthogonal
to the LQG-like native span.  The Schur correction therefore removes nothing
from the area sector, leaving the whole area currency as residual.  The
numerical values are computed in \texttt{simulate\_step2\_lqg\_xi.py}.
\end{proof}

\begin{remark}[Complementarity]
In the Step 1 toy, the holographic entanglement lens closes the area-shadow
sector and leaves the amplitude/P3 sector.  In the Step 2 toy, the LQG-like
lens closes the amplitude/P3 sector and leaves the area-shadow sector.  The
two imports are therefore complementary partial recognition sources on this
declared carrier.
\end{remark}

\begin{remark}[Mode C setup only]
Adding both source lenses closes the uncoupled finite toy.  Step 2 does not
promote this to the coupled E018 root because no operational recombination
bridge between the two source records has been audited.
\end{remark}

\end{document}
